\subsection{Train / Validation / Test Split}

The last 30 days of the series (Nov 1, 2023 -- Nov 30, 2023) are held
out as the test set; everything earlier is training data. Within the
training data, the last 60 days are reserved as a validation split,
used by the LSTM and PatchTST scripts for early stopping. Total
\textbf{training data: 1\,897 days}; \textbf{validation: 60 days
inside training}; \textbf{test: 30 days}. None of the test days are
seen by any model during training or hyper-parameter selection.

\subsection{Rolling-Origin Evaluation Protocol}

We use rolling-origin (also known as walk-forward) evaluation with
\textbf{four non-overlapping windows of seven days each}. At window
$i$ the model sees the full training set plus the first
$i \cdot 7$ days of the test set as history, and produces a 7-day
forecast for days $i \cdot 7 + 1, \dots, (i+1) \cdot 7$ of the test
set. Metrics are computed on each window separately and then averaged.
This protocol mimics what a real operator would do: every Monday they
have access to the previous week's actuals and need to forecast the
upcoming week. With 30 test days we get four such weekly windows; the
remaining two days are dropped (we keep windows non-overlapping so
metrics across windows are independent).

The classical models (SARIMA family, ARIMAX, ARMA, ARIMA, AutoARIMA)
are \emph{refit at every origin} so the model has access to the
freshest history. The LSTM is trained once on data prior to the test
window and rolled forward autoregressively. The PatchTST is also
trained once before the test window and uses a direct multi-step head
(no autoregression) --- at each origin it consumes the last 56 days
of history and outputs the 7-day forecast in one shot. Foundation
models (TimesFM, Chronos) operate zero-shot: they are not trained at
all, just queried with the history at each origin.

\subsection{Metrics}

We report seven metrics per window and aggregate them as
mean~$\pm$~standard deviation across the four windows:
\begin{itemize}[leftmargin=*]
    \item \textbf{MAE} (mean absolute error) in kWh --- our primary
    metric, because it is in the same units as the target and easy to
    interpret operationally (``the model is off by 25 kWh on
    average''). The standard deviation across windows is at least as
    important as the mean: a model with low mean but high variance is
    risky in production.
    \item \textbf{RMSE} (root mean squared error) in kWh, which
    penalises large errors more than MAE.
    \item \textbf{MAPE} and \textbf{sMAPE} (mean absolute percentage
    error and its symmetric variant). MAPE blows up on zero actuals
    so we report both.
    \item \textbf{MASE} (mean absolute scaled error), scaled by the
    in-sample MAE of the seasonal-naive forecast. MASE $< 1$ means
    the model beats the trivial seasonal baseline; MASE $= 0.65$
    means it cuts that error by 35\%.
    \item \textbf{Bias} (mean signed error) in kWh --- positive means
    the model over-predicts on average. Useful for spotting
    systematic offsets.
    \item \textbf{R\textsuperscript{2}} on the test window. R\textsuperscript{2}
    can be negative when the model is worse than the mean of the
    test window, which is common with very short windows; we still
    report it for completeness.
\end{itemize}
Throughout the paper we rank models by MAE because it is the most
operationally meaningful metric. The full table with every metric is
in Table~\ref{tab:leaderboard} of Section~\ref{sec:results}.

\subsection{Reproducibility and Heavy-Training Scripts}

All the heavy training lives in standalone Python scripts under
\texttt{scripts/}:
\texttt{sarima\_grid\_search.py} for the AIC grid,
\texttt{lstm\_train.py} for the LSTM sweep, and
\texttt{patchtst\_train.py} for the PatchTST seed ensemble. The
notebooks load the resulting cached predictions and apply the shared
\texttt{evaluation.py} rolling-origin routine, so a single notebook
never has to retrain a transformer to produce its leaderboard row.
Training the PatchTST seed ensemble takes $\approx 2$~minutes per
seed on an RTX 5070~Ti; the full pipeline (every model, every metric)
runs end-to-end in under ten minutes from a cold checkout.
